\subsection{Přímky v prostoru}
\label{subsec:Lines_in_space}

Stejně jako v rovině je přímka určena jedním svým bodem a směrem,
kterým prochází. Tento směr můžeme vyjádřit pomocí nenulového
směrového vektoru. Známe-li tedy bod
\[
    A[x_A;y_A;z_A]
\]
a nenulový vektor
\[
    \vec{u}=(u_x;u_y;u_z),
\]
je tím určena právě jedna přímka procházející bodem $A$ ve směru
vektoru $\vec{u}$.

Přímku můžeme určit také pomocí dvou různých bodů $A$ a $B$.
Jejím směrovým vektorem potom může být například vektor
\[
    \overrightarrow{AB}=B-A.
\]
Stejně jako v rovině tedy platí, že dvěma různými body prochází právě
jedna přímka.

Analytický popis přímky v prostoru se však od popisu přímky v rovině
v jednom důležitém ohledu liší. V rovině jsme mohli přímku popsat
jedinou obecnou rovnicí
\[
    ax+by+c=0.
\]
V prostoru již jediná lineární rovnice o proměnných $x$, $y$ a $z$
obecně nepopisuje přímku, ale rovinu. Přímku v prostoru proto budeme
nejprve popisovat pomocí \textbf{parametrického vyjádření}, které
přirozeně vychází z bodu přímky a jejího směrového vektoru.

Později si ukážeme, že přímku v prostoru lze popsat také jako
průsečnici dvou rovin.